\section{Discussion and Conclusion}
We introduced rotation-equivariant spherical neural fields and used them to create RENI++: a natural illumination prior. We demonstrated how random samples from RENI++ produce plausible illumination environments and used RENI++ for environment completion and inverse rendering. 

Our modifications in RENI++ have overcome two of the original limitations of the work, namely the \(O(n^{2})\) complexity of the Gram matrix and the reproduction of higher frequency details. However, some limitations remain. Because RENI++ has a prior for natural illumination, unlike SH, the performance of RENI++ decreases when fitted to indoor scenes. SH can also handle unnatural illuminations, such as a scene with multiple suns, something that will not lie within the latent space of RENI++ since it saw no such examples at training time. 

Human vision has complex interactions between illumination, geometry and texture priors. For example, the Hollow Face Illusion \cite{hill_independent_1993} arises from face geometry priors overriding the lighting from above illumination prior; while in the Bas relief ambiguity \cite{belhumeur_bas-relief_1999}, geometric priors cause incorrect lighting estimation. Our model discounts these interactions, learning an illumination prior independently but not its interactions with other cues. 

RENI++ is an auto-decoder based generative model. Whilst powerful, these suffer from a number of limitations. The dimensionality of their latent spaces, and consequently the number of optimisable parameters, increases proportionally with the dataset size. Since the latent space is randomly initialised and each latent code optimised independently, the latent space is often poorly structured resulting in poor quality generations and a more challenging optimisation landscape for inverse tasks. To mitigate these issues, future research could explore using a vector neuron transformer encoder. Additionally, gradient origin networks could provide a more structured and well-behaved latent representation, facilitating improved generation and usefulness as a prior in inverse tasks.

We strictly allow only $O(2)$ equivariance. However, considering typical camera coordinate systems, the up axis will sometimes not align with gravity when the camera is pointed up or down. For inverse problems, this would mean the gravity vector would need to be explicitly estimated (an accelerometer would resolve this). Alternatively, we could build our model with full $O(3)$ equivariance but then learn a prior over the space of camera poses relative to gravity. Finally, there are many opportunities to deploy RENI++ in larger inverse rendering pipelines where it could be a simple drop-in replacement for SH while another potential application is to use RENI++ for LDR to HDR panoramic image reconstruction. 